% ============================================================
%  SCHOLA DOMESTICA — Magister Mathematica
%  Level 1 | Week 13 | Day 2
%  Topic: Fact Families — Addition and Subtraction Partners from 8
% ============================================================
\documentclass[12pt, letterpaper]{article}
\usepackage{sd_style}

\renewcommand{\lessonweek}{13}
\renewcommand{\lessonnum}{2}
\renewcommand{\lessontitle}{Fact Families: Addition and Subtraction Partners}

% IMAGE PROMPT (Beatrix Potter watercolour style):
% Three little squirrels on an oak branch — one brown, one grey, one russet —
% looking curiously at two acorns on a mossy stump below. Autumn light. White background, no text.

\begin{document}

\lessonheader{Week 13}{2}{Fact Families: Addition and Subtraction Partners}

{\small Name:\ansblanklg \hfill Date:\ansblank}
\goldrule

\begin{instructbox}
A \textbf{fact family} uses the same three numbers to make \textbf{four} sentences — two addition, two subtraction.
\begin{center}
\begin{tikzpicture}
  \node[draw=sd-blue,circle,thick,minimum size=0.9cm,fill=sd-light](top) at (0,1.6){\textbf{\large 8}};
  \node[draw=sd-blue,circle,thick,minimum size=0.75cm,fill=sd-light](bl) at (-1.4,0){\textbf{3}};
  \node[draw=sd-blue,circle,thick,minimum size=0.75cm,fill=sd-light](br) at (1.4,0){\textbf{5}};
  \draw[thick,sd-blue](top)--(bl); \draw[thick,sd-blue](top)--(br); \draw[thick,sd-gold](bl)--(br);
\end{tikzpicture}
\end{center}
Numbers \textbf{3, 5, 8}:\quad $3+5=8$ \quad $5+3=8$ \quad $8-5=3$ \quad $8-3=5$
\end{instructbox}

\bigskip
\noindent{\color{sd-blue}\large\textbf{Part I — Complete Each Fact Family}} \hfill \textit{(Problems 1–8)}

\medskip\noindent\textit{Use the three numbers to write all four sentences.}

\bigskip
\noindent\textbf{Family 1 — Numbers: 2, 6, 8}

\medskip
\noindent\begin{tabular}{@{}p{5cm}p{5cm}p{5cm}@{}}
\prob $2+6=\ansblank$ & \prob $6+2=\ansblank$ & \\[12pt]
\prob $8-6=\ansblank$ & \prob $8-2=\ansblank$ &
\end{tabular}

\bigskip
\noindent\textbf{Family 2 — Numbers: 3, 4, 7}

\medskip
\noindent\begin{tabular}{@{}p{5cm}p{5cm}p{5cm}@{}}
\prob $3+4=\ansblank$ & \prob $4+3=\ansblank$ & \\[12pt]
\prob $7-4=\ansblank$ & \prob $7-3=\ansblank$ &
\end{tabular}

\bigskip\bigskip
\noindent{\color{sd-blue}\large\textbf{Part II — Find the Missing Number}} \hfill \textit{(Problems 9–14)}

\medskip\noindent\textit{Think of the fact family. Write the missing number.}

\bigskip
\noindent\begin{tabular}{@{}p{5cm}p{5cm}p{5cm}@{}}
\prob $5+\ansblank=8$ & \prob $\ansblank+4=7$ & \prob $8-\ansblank=4$ \\[14pt]
\prob $7-\ansblank=3$ & \prob $\ansblank-5=3$ & \prob $\ansblank-4=2$
\end{tabular}

\bigskip\bigskip
\noindent{\color{sd-blue}\large\textbf{Part III — Build Your Own Triangle}} \hfill \textit{(Problems 15–18)}

\medskip\noindent\textit{Add the two small numbers to find the top. Then write two sentences.}

\bigskip
\noindent\begin{tabular}{@{}p{7.5cm}p{7.5cm}@{}}
\textbf{A.}
\begin{tikzpicture}
  \node[draw=sd-blue,circle,thick,minimum size=0.9cm,fill=sd-light](top) at (0,1.4){\ansblank};
  \node[draw=sd-blue,circle,thick,minimum size=0.75cm,fill=sd-light](bl) at (-1.2,0){\textbf{1}};
  \node[draw=sd-blue,circle,thick,minimum size=0.75cm,fill=sd-light](br) at (1.2,0){\textbf{7}};
  \draw[thick,sd-blue](top)--(bl);\draw[thick,sd-blue](top)--(br);\draw[thick,sd-gold](bl)--(br);
\end{tikzpicture}
&
\textbf{B.}
\begin{tikzpicture}
  \node[draw=sd-blue,circle,thick,minimum size=0.9cm,fill=sd-light](top) at (0,1.4){\ansblank};
  \node[draw=sd-blue,circle,thick,minimum size=0.75cm,fill=sd-light](bl) at (-1.2,0){\textbf{2}};
  \node[draw=sd-blue,circle,thick,minimum size=0.75cm,fill=sd-light](br) at (1.2,0){\textbf{5}};
  \draw[thick,sd-blue](top)--(bl);\draw[thick,sd-blue](top)--(br);\draw[thick,sd-gold](bl)--(br);
\end{tikzpicture}
\end{tabular}

\medskip
\prob \ansblank $+$ \ansblank $=$ \ansblank \hfill \prob \ansblank $-$ \ansblank $=$ \ansblank

\medskip
\prob \ansblank $+$ \ansblank $=$ \ansblank \hfill \prob \ansblank $-$ \ansblank $=$ \ansblank

\bigskip
\begin{checkbox}
\textbf{\color{sd-green}Three Numbers, Four Friends!}\\[3pt]
Every fact family has one \textbf{big number} (the whole) and two \textbf{small numbers} (the parts).\\
The big number is always on top. The parts go on the bottom corners.\\[3pt]
\textit{How many fact families have 6 as the big number?} \hfill \textit{(Hint: 3 families!)}
\end{checkbox}

\end{document}
